\section{Ôn tập học kỳ 2 - Đề số 2}
\Opensolutionfile{ans}[ans/ans-2-OTHK2-De2]
\caulc

\begin{ex}%[Nguyễn Tuấn, dự án sáng tác đề 12]%[2D4N1-4]
	Họ nguyên hàm của hàm số $f(x)=\mathrm{e}^x + x$ là
	\choice
	{$\mathrm{e}^x + x^2 + C$}
	{\True $\mathrm{e}^x + \dfrac{1}{2}x^2 + C$}
	{$\dfrac{1}{x  + 1}\mathrm{e}^x + \dfrac{1}{2}x^2 + C$}
	{$\mathrm{e}^x + 1 + C$}
	\loigiai
	{
		Ta có
		$$\displaystyle\int f(x) \mathrm{\,d}x = \displaystyle\int(\mathrm{e}^x+x)\mathrm{\,d}x = \displaystyle\int \mathrm{e}^x \mathrm{\,d}x + \displaystyle\int x \mathrm{\,d}x = \mathrm{e}^x+\dfrac{1}{2}x^2+C, \text{ với } C \text{ là hằng số}.$$
	}
\end{ex}
\begin{ex}%[Nguyễn Tuấn, dự án sáng tác đề 12]%[2D4H1-2]
	Cho hàm số $y = \dfrac{2x^4 + 3}{x^2}$. Khẳng định nào sau đây là đúng?
	\choice
	{$ \displaystyle\int\limits f(x)\mathrm{d}x = \dfrac{2x^3}{3} + \dfrac{3}{2x} + \mathrm{C} $}
	{\True $ \displaystyle\int\limits f(x)\mathrm{d}x = \dfrac{2x^3}{3} - \dfrac{3}{x} + \mathrm{C} $}
	{$ \displaystyle\int\limits f(x)\mathrm{d}x = \dfrac{2x^3}{3} + \dfrac{3}{x} + \mathrm{C} $}
	{$ \displaystyle\int\limits f(x)\mathrm{d}x = 2x^3 - \dfrac{3}{x} + \mathrm{C} $}
	\loigiai{
		Ta có $ \displaystyle\int\limits f(x)\mathrm{d}x = \displaystyle\int\limits \dfrac{2x^4 + 3}{x^2}\mathrm{d}x = \displaystyle\int\limits \left(2x^2 + \dfrac{3}{x^2}\right)\mathrm{d}x = \dfrac{2x^3}{3} - \dfrac{3}{x} + \mathrm{C} $.
	}
\end{ex}
\begin{ex}%[Nguyễn Tuấn, dự án sáng tác đề 12]%[2D4N2-3]
	Giá trị của $\displaystyle\int\limits_0^{\frac{\pi}{2}} \cos x\mathrm{\,d}x$ bằng
	\choice
	{$0$}
	{\True $1$}
	{$\dfrac{\pi}{2}$}
	{$\pi$}
	\loigiai
	{
		Ta có $\displaystyle\int\limits_0^{\frac{\pi}{2}} \cos x\mathrm{\,d}x = \sin x\Big|_0^{\frac{\pi}{2}} = \sin\dfrac{\pi}{2}-\sin 0 = 1$.
	}
\end{ex}
\begin{ex}%[Nguyễn Tuấn, dự án sáng tác đề 12]%[2D4H2-1]
	Tích phân $ \displaystyle\int\limits_{1}^{2}[4f(x) - 2x]\mathrm{d}x = 1$. Khi đó $ \displaystyle\int\limits_{1}^{2}f(x)\mathrm{d}x$ bằng
	\choice
	{\True $ 1$}
	{$ - 3 $}
	{$ 3 $}
	{$ -1$}
	\loigiai{
		Ta có:
		{\allowdisplaybreaks
			\begin{eqnarray*}
				& & \displaystyle\int\limits_{1}^{2}[4f(x) - 2x]\mathrm{d}x = 1\\
				&\Leftrightarrow & 4\displaystyle\int\limits_{1}^{2}f(x)\mathrm{d}x-2\displaystyle\int\limits_{1}^{2}x\mathrm{d}x = 1 \\
				&\Leftrightarrow & 4\displaystyle\int\limits_{1}^{2}f(x)\mathrm{d}x - 2 \cdot \dfrac{x^2}{2} \Big|_1^2 = 1\\
				&\Leftrightarrow & 4\displaystyle\int\limits_{1}^{2}f(x)\mathrm{d}x = 4\\
				&\Leftrightarrow & \displaystyle\int\limits_{1}^{2}f(x)\mathrm{d}x = 1.
		\end{eqnarray*}}
	}
\end{ex}
\begin{ex}%[Nguyễn Tuấn, dự án sáng tác đề 12]%[2D4N3-1]
	\immini{Diện tích phần hình phẳng gạch chéo trong hình vẽ bên được tính theo công thức nào dưới đây?
		\choice
		{$\displaystyle\int\limits_{-1}^2\left(2x^2-2x-4\right)\mathrm{\,d}x$}
		{$\displaystyle\int\limits_{-1}^2(-2x+2)\mathrm{\,d}x$}
		{$\displaystyle\int\limits_{-1}^2(2x-2)\mathrm{\,d}x$}
		{\True $\displaystyle\int\limits_{-1}^2\left(-2x^2+2x+4\right)\mathrm{\,d}x$}}
	{\begin{tikzpicture}[scale=.7, font=\footnotesize, line join=round, line cap=round, >=stealth]
			\draw[->] (-1.5,0) -- (3,0)node[below]{\footnotesize $x$};
			\draw (-1,0) circle (.5pt)node[below]{\footnotesize $-1$};
			\draw (2,0) circle (.5pt)node[above]{\footnotesize $2$};
			\draw[->,color=black] (0,-2.5) -- (0,3.5)node[below left]{\footnotesize $y$};
			\fill[pattern=north west lines] plot[smooth,samples=100,domain=-1:2] (\x,{(\x)^2-2*(\x)-1})--plot[smooth,samples=100,domain=2:-1] (\x,{-(\x)^2+3});
			\draw[thick,smooth,samples=100,domain=-1.5:2.2] plot(\x,{-(\x)^2+3});
			\draw[thick,smooth,samples=100,domain=-1.2:3] plot(\x,{(\x)^2-2*(\x)-1});
			\draw[dashed] (2,0) -- (2,-1) (-1,0) -- (-1,2);
			\filldraw[fill=white] (0,0) circle (1pt)node[shift={(-45:6pt)}]{\footnotesize $O$};
			\draw (2,3) node{\footnotesize $y=-x^2+3$};
			\draw (-1,-2.2) node{\footnotesize $y=x^2-2x-1$};
	\end{tikzpicture}}
	\loigiai{
		$S=\displaystyle\int\limits_{-1}^2\left[\left(-x^2+3\right)-\left(x^2-2x-1\right)\right]\mathrm{\,d}x=\displaystyle\int\limits_{-1}^2\left(-2x^2+2x+4\right)\mathrm{\,d}x$.
	}
\end{ex}
\begin{ex}%[Nguyễn Tuấn, dự án sáng tác đề 12]%[2D4H3-3]
	Thể tích của vật thể tròn xoay thu được khi quay hình phẳng (phần gạch sọc của hình vẽ) xung quanh trục $Ox$ bằng
\begin{center}
	\begin{tikzpicture}[line join=round, line cap = round, >=stealth, scale=.8,font=\footnotesize,transform shape]
		\draw[->] (0,-.5)--(0,4) node[above left]{$y$};
		\draw[->] (-.5,0)--(7,0) node[above]{$x$};
		\def \f(#1){sqrt(6*#1 - #1*#1)};
		\draw[color=black,smooth,samples=100,domain=0:6] plot(\x,{\f(\x)});
		\draw 
		(6.5,2) node{$y=\sqrt{6x-x^2}$}
		(-.5,-.5)--(3.5,3.5) node[right]{$y=x$}
		;
		\draw[dashed](3,-.5)--(3,3.5);
		\fill[pattern=north west lines] plot[domain=3:6](\x,{\f(\x)})--(0,0)--cycle;
		\foreach \x/\y/\z/\g in {0/0/O/-45,3/0/3/-45,6/0/6/-90}
		\fill[black] (\x,\y) circle(1pt) ($(\x,\y)+(\g:3mm)$) node{$\z$};
	\end{tikzpicture}
\end{center}
	\choice
	{$24\pi$}
	{\True $27\pi$}
	{$25\pi$}
	{$26\pi$}
	\loigiai{
		\begin{center}
			\begin{tikzpicture}[line join=round, line cap = round, >=stealth, scale=.8,font=\footnotesize,transform shape]
				\draw[->] (0,-.5)--(0,4) node[above left]{$y$};
				\draw[->] (-.5,0)--(7,0) node[above]{$x$};
				\def \f(#1){sqrt(6*#1 - #1*#1)};
				\draw[color=black,smooth,samples=100,domain=0:6] plot(\x,{\f(\x)});
				\draw 
				(6.5,2) node{$y=\sqrt{6x-x^2}$}
				(-.5,-.5)--(3.5,3.5) node[right]{$y=x$}
				;
				\draw[dashed](3,-.5)--(3,3.5);
				\fill[pattern=north west lines] plot[domain=3:6](\x,{\f(\x)})--(0,0)--cycle;
				\foreach \x/\y/\z/\g in {0/0/O/-45,3/0/3/-45,6/0/6/-90}
				\fill[black] (\x,\y) circle(1pt) ($(\x,\y)+(\g:3mm)$) node{$\z$};
			\end{tikzpicture}
		\end{center}
		Chia vật thể bởi mặt phẳng $x=3$ thì vật thể chia thành $2$ phần:
		\begin{itemize}
			\item Phần $1$ là khối tròn xoay khi quay hình phẳng giới hạn bởi các đường $y=x$, trục hoành và hai đường thẳng $x=0$, $x=3$ quanh trục $Ox$. Phần này có thể tích là:  
			$$V_1=\pi\displaystyle\int\limits_{0}^{3}x^2\mathrm{\,d}x=\pi\displaystyle\int\limits_{0}^{3}x^2\mathrm{\,d}x = \pi \dfrac{x^3}{3}\Big|_0^{3}=9\pi.$$
			\item Phần $2$ là khối tròn xoay khi quay hình phẳng giới hạn bởi các đường $y=\sqrt{6x-x^2}$, trục hoành và hai đường thẳng $x=3$, $x=6$ quanh trục $Ox$. Phần này có thể tích là:
			$$V_2=\pi\displaystyle\int\limits_{3}^{6}\left(\sqrt{6x - x^2} \right)^2\mathrm{\,d}x=\pi\displaystyle\int\limits_{3}^{6}\left(6x - x^2 \right)\mathrm{\,d}x = \pi\left( 3x^2 - \dfrac{x^3}{3} \right)\Big|_3^{6}=18\pi.$$
		\end{itemize}
		Vậy thể tích vật thể là $V=V_1+V_2=27\pi$.
	}
\end{ex}
\begin{ex}%[Nguyễn Tuấn, dự án sáng tác đề 12]%[2H5N1-2]
	Trong không gian $Oxyz$, cho mặt phẳng $(P)\colon 3x-z+2=0$. Véc-tơ nào dưới đây là một véc-tơ pháp tuyến của $(P)$?
	\choice
	{$\overrightarrow{n}=\left(-1;0;-1\right)$}
	{$\overrightarrow{n}=\left(3;-1;2\right)$}
	{$\overrightarrow{n}=\left(3;-1;0\right)$}
	{\True $\overrightarrow{n}=\left(3;0;-1\right)$}
	\loigiai{
		Mặt phẳng $(P)\colon 3x-z+2=0$ có một véc-tơ pháp tuyến là $\overrightarrow{n}=\left(3;0;-1\right)$.
	}
\end{ex}
\begin{ex}%[Nguyễn Tuấn, dự án sáng tác đề 12]%[2H5V2-3]
	Trong không gian $Oxyz$, cho điểm $A(1 ; 3 ; 2)$, mặt phẳng $(P)\colon 2 x-y+z-10=0$  và đường thẳng  $d\colon \dfrac{x+2}{2}=\dfrac{y-1}{1}=\dfrac{z-1}{-1}$. Đường thẳng $\Delta$ cắt $(P)$ và $d$ lần lượt tại hai điểm $M, N$ sao cho $A$ là trung điểm của đoạn $MN$. Biết rằng $\vec{u}=(a ; b ; 1)$ là một vectơ chỉ phương của $\Delta$, giá trị của $a+b$ bằng
	\choice
	{$11$}
	{\True $-11$}
	{$3$}
	{$-3$}
	\loigiai{
		Vì $N\in d$ nên $N\left(-2+2t;1+t;1-t \right)$ $t\in\mathbb{R}$. \\ Mặt khác, $A$ là trung điểm đoạn $MN$ nên $M(4-2t;5-t;3+t)$. \\
		Do $M \in (P)$ nên $2(4-2t)-(5-t)+3+t-10=0 \Leftrightarrow t = -2$. \\ Suy ra $M(8;7;1)$ và $N(-6;-1;3) \Rightarrow \vec{MN} =(-14;-8;2)$. \\
		Khi đó, đường thẳng $\Delta$ có một véc-tơ chỉ phương là  $\vec{u}=(-7 ; -4 ; 1)$. \\
		Suy ra $a = -7$ và $b = -4$. Vậy $a + b = -11$.
	}
\end{ex}
\begin{ex}%[Nguyễn Tuấn, dự án sáng tác đề 12]%[2H5H2-7]
	Trong không gian $Oxyz$, cô-sin của góc giữa đường thẳng chứa trục $Oy$ và mặt phẳng $(P)\colon 4x-3y+\sqrt{2}z-7=0$ bằng
	\choice
	{\True $\dfrac{\sqrt{2}}{\sqrt{3}}$}
	{$\dfrac{2}{\sqrt{3}}$}
	{$\dfrac{1}{\sqrt{3}}$}
	{$\dfrac{4}{\sqrt{3}}$}
	\loigiai{
		Véc-tơ chỉ phương của đường thẳng chứa trục $Oy$ là $\overrightarrow{j}=(0;1;0)$ và véc-tơ pháp tuyến của $(P)$ là $\overrightarrow{n}\left(4;-3;\sqrt{2}\right)$. Gọi $\alpha$ là góc giữa đường thẳng và mặt phẳng trên thì
		\[ \sin \alpha =\left| \cos \left(\overrightarrow{j},\overrightarrow{n}\right) \right| =\frac{1}{\sqrt{3}}. \]
		Suy ra $\cos \alpha =\sqrt{1-\sin^2\alpha}=\dfrac{\sqrt{2}}{\sqrt{3}}$.
	}
\end{ex}
\begin{ex}%[Nguyễn Tuấn, dự án sáng tác đề 12]%[2H5N3-2]
	Trong không gian $Oxyz$, mặt cầu $(S)\colon (x-1)^2+(y-2)^2+(z+3)^2=4$ có bán kính bằng
	\choice
	{$4$}
	{\True $2$}
	{$\sqrt{2}$}
	{$16$}
	\loigiai{Mặt cầu $(S)$ có bán kính $R=2$.}
\end{ex}
\begin{ex}%[Nguyễn Tuấn, dự án sáng tác đề 12]%[2D5H2-2]
	Cho $\mathrm{P}(A)=\dfrac{2}{5}$; $\mathrm{P}\left( B\mid A\right)=\dfrac{1}{3}$. Giá trị của $\mathrm{P}(AB)$ là 
	\choice
	{\True $\dfrac{2}{15} $}
	{$ \dfrac{3}{16}$}
	{$ \dfrac{1}{5}$}
	{$ \dfrac{4}{15}$}
	\loigiai{
		$\mathrm{P}(AB)=\mathrm{P}(A)\cdot \mathrm{P}\left( B\mid A\right)=\dfrac{2}{5}\cdot \dfrac{1}{3}=\dfrac{2}{15}$.
	}
\end{ex}
\begin{ex}%[Nguyễn Tuấn, dự án sáng tác đề 12]%[2D5V2-3]
	An có một túi gồm một số chiếc kẹo cùng loại, chỉ khác màu, trong đó có $6$ chiếc kẹo sô-cô-la đen, còn lại là $4$ chiếc kẹo sô-cô-la trắng. An lấy ngẫu nhiên $1$ chiếc kẹo trong túi để cho Bình, rồi lại lấy ngẫu nhiên tiếp 1 chiếc kẹo nữa trong túi và cũng đưa cho Bình. Xác suất để Bình nhận được $2$ chiếc kẹo sô-cô-la đen là 
	\choice
	{\True $\dfrac{1}{3} $}
	{$ \dfrac{1}{4}$}
	{$ \dfrac{2}{5}$}
	{$ \dfrac{3}{7}$}
	\loigiai{
		Gọi $A$ là biến cố: \lq\lq An lấy lần 1 được 1 chiếc kẹo sô-cô-la đen\rq\rq\, và $B$ là biến cố: \lq\lq An lấy lần 2 được 1 chiếc kẹo sô-cô-la đen\rq\rq.\\
		Khi đó $AB$ là biến cố: \lq\lq Cả hai lần đều lấy được kẹo sô-cô-la đen\rq\rq.\\
		Ta có $\mathrm{P}(A)=\dfrac{n(A)}{n(\Omega)}=\dfrac{6}{10}$.\\
		Sau khi lấy 1 chiếc kẹo sô-cô-la đen thì xác suất để chọn 1 chiếc kẹo sô-cô-la đen trong hộp đựng $5$ chiếc kẹo sô-cô-la đen, còn lại là $4$ chiếc kẹo sô-cô-la trắng là $\mathrm{P}(B\mid A)=\dfrac{5}{9}$.\\
		Khi đó $\mathrm{P}(AB)=\mathrm{P}(A)\cdot \mathrm{P}(B\mid A)=\dfrac{6}{10}\cdot \dfrac{5}{9}=\dfrac{1}{3}$.\\
		Xác suất để Bình nhận được $2$ chiếc kẹo sô-cô-la đen là $\dfrac{1}{3}$.
	}
\end{ex}
\Closesolutionfile{ans}
\indapan{6}{ans/ans-2-OTHK2-De2}


\cauds
\Opensolutionfile{ans}[ans/ans-2-OTHK2]
\begin{ex}%Cau1D %[2H5N3-3]
Cho $\left(S\right): x^2+y^2+z^2-2x-4y+6z-67=0$. Các mệnh đề sau đúng hay sai ?
\choiceTF
	{Mặt cầu $\left( S\right)$ có tâm $I\left(1;2;3\right)$}
	{\True Bán kính mặt cầu $\left( S\right)$ là $R=9$}
	{Cho mặt phẳng $\left(P\right): 2x-2y+z-13=0$. Khi đó $\left( P\right)$ tiếp xúc với $\left(S\right)$}
	{\True Cho đường thẳng $\left(\Delta\right) :\heva{&x=1+t\\&y=2\\&z=-4+7t}$. Khi đó $\left(\Delta\right)$ và $\left(S\right)$ cắt nhau tại hai điểm}
	\loigiai{
		\begin{itemchoice}
			\itemch Sai. Mặt cầu $\left( S\right)$ có tâm $I\left(1;2;-3\right)$.
			\itemch Đúng. Ta có $R=\sqrt{1^2+2^2+\left(-3\right)^2+67}=\sqrt{81}=9$
			\itemch Sai. Vì $\mathrm{d}\left(I,\left(P\right)\right)=\dfrac{\vert 2\cdot1-2\cdot2+\left(-3\right)-13\vert}{\sqrt{2^2+2^2+\left(-1\right)^2}}=6< R=9$. Suy ra $\left(P\right)$ cắt $\left(S\right)$ theo giao tuyến là một đường tròn.
			\itemch Đúng. Tọa độ giao điểm là nghiệm của hệ phương trình\\
	 $\heva{&x=1+t\\&y=2\\&z=-4+7t\\&x^2+y^2+z^2-2x-4y+6z-67=0}\Rightarrow  \hoac{&t=0\\&t=1}$.\\
	 Với $t=0$, thay $t=0$ vào $\left(\Delta\right)$ ta được tọa độ $A\left(1;2;-4\right)$\\
	 Với $t=1$, thay $t=1$ vào $\left(\Delta\right)$ ta được tọa độ $B\left(2;2;3\right)$.\\
	 Vậy, $\left(\Delta\right)$ và $\left(S\right)$ cắt nhau tại hai điểm $A$ và $B$.
		\end{itemchoice}
	}
\end{ex}

\begin{ex}%Cau2D %[2D5V1-2]
Trong một hộp có 20 viên bi xanh và 4 viên bi đỏ, các viên bi đều có hình dạng và kích thước giống nhau. Một học sinh lấy ngẫu nhiên lần lượt $2$ viên bi (lấy không hoàn lại) trong hộp.
Các khẳng định nào dưới đây đúng hay sai ? 
\choiceTF
	{Xác suất để lần thứ nhất lấy được viên bi đỏ là $\dfrac{1}{5}$}
	{\True Xác suất để lần thứ hai lấy được viên bi đỏ, biết lần thứ nhất lấy được viên bi đỏ, là $\dfrac{3}{23}$}
	{\True Xác suất để cả hai lần đều lấy được viên bi đỏ là $\dfrac{1}{46}$}
	{\True Xác suất để ít nhất một lần lấy được viên bi xanh là $\dfrac{45}{46}$}
	\loigiai{
		\begin{itemchoice}
			\itemch Sai. Xét biến cố:\\
	$A$:"Lần thứ nhất lấy được viên bi đỏ". Khi đó xác suất lần đầu lấy được viên bi đỏ là: $P(A)=\dfrac{4}{24}=\dfrac{1}{6}$.
			\itemch Đúng. Khi lần thứ nhất đã lấy được viên bi đỏ thì trong hộp còn $3$ viên bi đỏ. Gọi biến cố $B$: "Lần thứ hai lấy được viên bi đỏ", khi đó xác suất lần thứ hai lấy được viên bi đỏ khi biết lần thứ nhất lấy được viên bi đỏ là: $P(B \vert A)=\dfrac{3}{23}$
			\itemch Đúng. Xác suất để cả hai lần đều lấy được viên bi đỏ là: $P(C)=P(A \cap B)=P(A) \cdot P(B \vert A) = \dfrac{1}{6} \cdot \dfrac{3}{23}=\dfrac{1}{46}$. 
			\itemch Đúng. Xác suất để có ít nhất một lần lấy được viên bi xanh là:\\
			 $P(\overline{C})=1-P(C)= 1-\dfrac{1}{46}=\dfrac{45}{46}$. 
		\end{itemchoice}
	}
\end{ex}

\begin{ex}%Cau3D %[2H5H2-7]
	Trong không gian với hệ tọa độ $Oxyz$, cho mặt phẳng $(P)\colon 3x-y+2z+2=0$ và đường thẳng $d\colon\dfrac{x}{3}=\dfrac{y-1}{2}=\dfrac{z+1}{-1}$. Xác định tính đúng, sai của mỗi mệnh đề sau
	\choiceTF
	{$(3;-1; -2)$ là véc-tơ pháp tuyến của $(P)$}
	{\True Góc giữa $(P)$ và $(d)$ là một góc nhọn}
	{\True Góc giữa $(P)$ và $(d)$ là $20^{\circ}$ (làm tròn đến hàng đơn vị của độ)}
	{Góc giữa trục $Ox$ và $d$ (làm tròn đến hàng đơn vị của độ) bằng $42^\circ$}
	\loigiai{
	Ta có 	$\overrightarrow{n}_{(P)}=(1;-1; 2)$ là một véc-tơ pháp tuyến của $(P)$ và $\overrightarrow{u}_d=(1; 2;-1)$ là một véc-tơ chỉ phương của $d$.
		\begin{itemchoice}
		\itemch Sai. Ta có $\overrightarrow{n}_{(P)}=(3;-1; 2)$.
		\itemch Đúng. 
		$\sin\left(d,(P)\right)=\dfrac{\left|\overrightarrow{n}_{(P)}\cdot\overrightarrow{u}_d\right|}{\left|\overrightarrow{n}_{(P)}\right|\cdot\left|\overrightarrow{u}_d\right|} =\dfrac{|9-2-2|}{\sqrt{14}\cdot\sqrt{14}}=\dfrac{5}{14}$.  \\
		Suy ra 	$0<\sin\left(d,(P)\right)\ne 1$ do đó góc giữa $(P)$ và $(d)$ là một góc nhọn.
			\itemch Sai. 
			Vậy $\sin\left(d,(P)\right)=\dfrac{\left|\overrightarrow{n}_{(P)}\cdot\overrightarrow{u}_d\right|}{\left|\overrightarrow{n}_{(P)}\right|\cdot\left|\overrightarrow{u}_d\right|} =\dfrac{|9-2-2|}{\sqrt{14}\cdot\sqrt{14}}=\dfrac{5}{14}\Rightarrow \left(d,(P)\right)\approx 21^{\circ}$.
			\itemch Sai.
			Ta có  $\overrightarrow{u}_{Ox}=\overrightarrow{i}=(1; 0;0)$.\\
			$\cos\left(d,Ox\right)=\dfrac{\left|\overrightarrow{u}_{d}\cdot\overrightarrow{i}\right|}{\left|\overrightarrow{u}_{d}\right|\cdot\left|\overrightarrow{i}\right|} =\dfrac{|3-0-0|}{\sqrt{14}\cdot\sqrt{1}}=\dfrac{3}{\sqrt{14}} \Rightarrow (d,Ox)\approx 37^{\circ}$.
			\end{itemchoice}}
\end{ex}

\begin{ex}%Cau4D %[2D5N2-4]
Xét một lô giày chiến sĩ được sản xuất bởi $3$ nhà máy $1$, $2$, $3$ với tỉ lệ lần lượt là $20\%$, $30\%$ và $50\%$. Xác suất giày hỏng của các nhà máy lần lượt là $0,001$; $0,005$ và $0,006$. Các khẳng định dưới đây đúng hay sai?
\choiceTF
	{\True Lấy ngẫu nhiên một chiếc giày từ lô hàng. Xác suất để chiếc giày lấy ra là giày hỏng ở nhà máy $1$ là $0,002$ }
	{\True Lấy ngẫu nhiên một chiếc giày từ lô hàng. Xác suất để chiếc giày lấy ra bị hỏng là $0,0065$}
	{Lấy ngẫu nhiên một chiếc giày từ lô hàng. Xác suất để chiếc giày lấy ra không bị hỏng là $0,9925$}
	{Lấy ngẫu nhiên ra một chiếc giày kiểm tra thì được kết quả là chiếc giày bị hỏng. Xác suất để chiếc giày đó do nhà máy $3$ sản xuất là $0,5$ (làm tròn đến chữ số thập phân thứ nhất)}
	\loigiai{ Gọi $A$: "lấy được giày hỏng", $A_i$:"lấy được giày ở nhà máy $i$" $\left(i=1,2,3\right)$. 
		\begin{itemchoice}
			\itemch Đúng. Ta có $P\left(A_1 \right)=0,2$, và $P\left(A \vert A_1 \right)=0,001$. Khi đó xác suất để chiếc giày lấy ra là giày hỏng ở nhà máy $1$ là $0,2 \cdot 0,001=0,002$.
			\itemch Đúng. Ta có $\left\lbrace A_1,A_2,A_3 \right\rbrace$ là hệ đầy đủ. Áp dụng công thức toàn phần ta có:
			\begin{eqnarray*}
&P\left(A\right) &= P\left(A_1\right)\cdot P\left(A\vert A_1 \right)+P\left(A_1\right)\cdot P\left(A\vert A_1 \right)+P\left(A_1\right)\cdot P\left(A\vert A_1 \right)\\
& &= 0,2 \cdot 0,001+0,3 \cdot 0,005+0,5 \cdot 0,006\\
& &= 0,0065
			\end{eqnarray*}
			\itemch Sai. Ta có $P\left(\overline{A}\right)=1-P\left(A\right)=1-0,0065=0,9935$.
			\itemch Sai. Ta có:$$P\left(A_3\vert A \right)=\dfrac{P\left(A_3\right)\cdot P\left(A \vert A_3\right)}{P\left(A\right)}=\dfrac{0,5\cdot 0,006}{0,0065}=\dfrac{6}{13}\approx 0,5$$
			Vậy xác suất để chiếc giày đó do nhà máy $3$ sản xuất là $0,5$.
		\end{itemchoice}
	}
\end{ex}

\Closesolutionfile{ans}
\indapan{3}{ans/ans-2-OTHK2}


\cauds
\Opensolutionfile{ans}[ans/ans-2-OTHK2]
\begin{ex}%Cau1D %[2H5N3-3]
Cho $\left(S\right): x^2+y^2+z^2-2x-4y+6z-67=0$. Các mệnh đề sau đúng hay sai ?
\choiceTF
	{Mặt cầu $\left( S\right)$ có tâm $I\left(1;2;3\right)$}
	{\True Bán kính mặt cầu $\left( S\right)$ là $R=9$}
	{Cho mặt phẳng $\left(P\right): 2x-2y+z-13=0$. Khi đó $\left( P\right)$ tiếp xúc với $\left(S\right)$}
	{\True Cho đường thẳng $\left(\Delta\right) :\heva{&x=1+t\\&y=2\\&z=-4+7t}$. Khi đó $\left(\Delta\right)$ và $\left(S\right)$ cắt nhau tại hai điểm}
	\loigiai{
		\begin{itemchoice}
			\itemch Sai. Mặt cầu $\left( S\right)$ có tâm $I\left(1;2;-3\right)$.
			\itemch Đúng. Ta có $R=\sqrt{1^2+2^2+\left(-3\right)^2+67}=\sqrt{81}=9$
			\itemch Sai. Vì $\mathrm{d}\left(I,\left(P\right)\right)=\dfrac{\vert 2\cdot1-2\cdot2+\left(-3\right)-13\vert}{\sqrt{2^2+2^2+\left(-1\right)^2}}=6< R=9$. Suy ra $\left(P\right)$ cắt $\left(S\right)$ theo giao tuyến là một đường tròn.
			\itemch Đúng. Tọa độ giao điểm là nghiệm của hệ phương trình\\
	 $\heva{&x=1+t\\&y=2\\&z=-4+7t\\&x^2+y^2+z^2-2x-4y+6z-67=0}\Rightarrow  \hoac{&t=0\\&t=1}$.\\
	 Với $t=0$, thay $t=0$ vào $\left(\Delta\right)$ ta được tọa độ $A\left(1;2;-4\right)$\\
	 Với $t=1$, thay $t=1$ vào $\left(\Delta\right)$ ta được tọa độ $B\left(2;2;3\right)$.\\
	 Vậy, $\left(\Delta\right)$ và $\left(S\right)$ cắt nhau tại hai điểm $A$ và $B$.
		\end{itemchoice}
	}
\end{ex}

\begin{ex}%Cau2D %[2D5V1-2]
Trong một hộp có 20 viên bi xanh và 4 viên bi đỏ, các viên bi đều có hình dạng và kích thước giống nhau. Một học sinh lấy ngẫu nhiên lần lượt $2$ viên bi (lấy không hoàn lại) trong hộp.
Các khẳng định nào dưới đây đúng hay sai ? 
\choiceTF
	{Xác suất để lần thứ nhất lấy được viên bi đỏ là $\dfrac{1}{5}$}
	{\True Xác suất để lần thứ hai lấy được viên bi đỏ, biết lần thứ nhất lấy được viên bi đỏ, là $\dfrac{3}{23}$}
	{\True Xác suất để cả hai lần đều lấy được viên bi đỏ là $\dfrac{1}{46}$}
	{\True Xác suất để ít nhất một lần lấy được viên bi xanh là $\dfrac{45}{46}$}
	\loigiai{
		\begin{itemchoice}
			\itemch Sai. Xét biến cố:\\
	$A$:"Lần thứ nhất lấy được viên bi đỏ". Khi đó xác suất lần đầu lấy được viên bi đỏ là: $P(A)=\dfrac{4}{24}=\dfrac{1}{6}$.
			\itemch Đúng. Khi lần thứ nhất đã lấy được viên bi đỏ thì trong hộp còn $3$ viên bi đỏ. Gọi biến cố $B$: "Lần thứ hai lấy được viên bi đỏ", khi đó xác suất lần thứ hai lấy được viên bi đỏ khi biết lần thứ nhất lấy được viên bi đỏ là: $P(B \vert A)=\dfrac{3}{23}$
			\itemch Đúng. Xác suất để cả hai lần đều lấy được viên bi đỏ là: $P(C)=P(A \cap B)=P(A) \cdot P(B \vert A) = \dfrac{1}{6} \cdot \dfrac{3}{23}=\dfrac{1}{46}$. 
			\itemch Đúng. Xác suất để có ít nhất một lần lấy được viên bi xanh là:\\
			 $P(\overline{C})=1-P(C)= 1-\dfrac{1}{46}=\dfrac{45}{46}$. 
		\end{itemchoice}
	}
\end{ex}

\begin{ex}%Cau3D %[2H5H2-7]
	Trong không gian với hệ tọa độ $Oxyz$, cho mặt phẳng $(P)\colon 3x-y+2z+2=0$ và đường thẳng $d\colon\dfrac{x}{3}=\dfrac{y-1}{2}=\dfrac{z+1}{-1}$. Xác định tính đúng, sai của mỗi mệnh đề sau
	\choiceTF
	{$(3;-1; -2)$ là véc-tơ pháp tuyến của $(P)$}
	{\True Góc giữa $(P)$ và $(d)$ là một góc nhọn}
	{\True Góc giữa $(P)$ và $(d)$ là $20^{\circ}$ (làm tròn đến hàng đơn vị của độ)}
	{Góc giữa trục $Ox$ và $d$ (làm tròn đến hàng đơn vị của độ) bằng $42^\circ$}
	\loigiai{
	Ta có 	$\overrightarrow{n}_{(P)}=(1;-1; 2)$ là một véc-tơ pháp tuyến của $(P)$ và $\overrightarrow{u}_d=(1; 2;-1)$ là một véc-tơ chỉ phương của $d$.
		\begin{itemchoice}
		\itemch Sai. Ta có $\overrightarrow{n}_{(P)}=(3;-1; 2)$.
		\itemch Đúng. 
		$\sin\left(d,(P)\right)=\dfrac{\left|\overrightarrow{n}_{(P)}\cdot\overrightarrow{u}_d\right|}{\left|\overrightarrow{n}_{(P)}\right|\cdot\left|\overrightarrow{u}_d\right|} =\dfrac{|9-2-2|}{\sqrt{14}\cdot\sqrt{14}}=\dfrac{5}{14}$.  \\
		Suy ra 	$0<\sin\left(d,(P)\right)\ne 1$ do đó góc giữa $(P)$ và $(d)$ là một góc nhọn.
			\itemch Sai. 
			Vậy $\sin\left(d,(P)\right)=\dfrac{\left|\overrightarrow{n}_{(P)}\cdot\overrightarrow{u}_d\right|}{\left|\overrightarrow{n}_{(P)}\right|\cdot\left|\overrightarrow{u}_d\right|} =\dfrac{|9-2-2|}{\sqrt{14}\cdot\sqrt{14}}=\dfrac{5}{14}\Rightarrow \left(d,(P)\right)\approx 21^{\circ}$.
			\itemch Sai.
			Ta có  $\overrightarrow{u}_{Ox}=\overrightarrow{i}=(1; 0;0)$.\\
			$\cos\left(d,Ox\right)=\dfrac{\left|\overrightarrow{u}_{d}\cdot\overrightarrow{i}\right|}{\left|\overrightarrow{u}_{d}\right|\cdot\left|\overrightarrow{i}\right|} =\dfrac{|3-0-0|}{\sqrt{14}\cdot\sqrt{1}}=\dfrac{3}{\sqrt{14}} \Rightarrow (d,Ox)\approx 37^{\circ}$.
			\end{itemchoice}}
\end{ex}

\begin{ex}%Cau4D %[2D5N2-4]
Xét một lô giày chiến sĩ được sản xuất bởi $3$ nhà máy $1$, $2$, $3$ với tỉ lệ lần lượt là $20\%$, $30\%$ và $50\%$. Xác suất giày hỏng của các nhà máy lần lượt là $0,001$; $0,005$ và $0,006$. Các khẳng định dưới đây đúng hay sai?
\choiceTF
	{\True Lấy ngẫu nhiên một chiếc giày từ lô hàng. Xác suất để chiếc giày lấy ra là giày hỏng ở nhà máy $1$ là $0,002$ }
	{\True Lấy ngẫu nhiên một chiếc giày từ lô hàng. Xác suất để chiếc giày lấy ra bị hỏng là $0,0065$}
	{Lấy ngẫu nhiên một chiếc giày từ lô hàng. Xác suất để chiếc giày lấy ra không bị hỏng là $0,9925$}
	{Lấy ngẫu nhiên ra một chiếc giày kiểm tra thì được kết quả là chiếc giày bị hỏng. Xác suất để chiếc giày đó do nhà máy $3$ sản xuất là $0,5$ (làm tròn đến chữ số thập phân thứ nhất)}
	\loigiai{ Gọi $A$: "lấy được giày hỏng", $A_i$:"lấy được giày ở nhà máy $i$" $\left(i=1,2,3\right)$. 
		\begin{itemchoice}
			\itemch Đúng. Ta có $P\left(A_1 \right)=0,2$, và $P\left(A \vert A_1 \right)=0,001$. Khi đó xác suất để chiếc giày lấy ra là giày hỏng ở nhà máy $1$ là $0,2 \cdot 0,001=0,002$.
			\itemch Đúng. Ta có $\left\lbrace A_1,A_2,A_3 \right\rbrace$ là hệ đầy đủ. Áp dụng công thức toàn phần ta có:
			\begin{eqnarray*}
&P\left(A\right) &= P\left(A_1\right)\cdot P\left(A\vert A_1 \right)+P\left(A_1\right)\cdot P\left(A\vert A_1 \right)+P\left(A_1\right)\cdot P\left(A\vert A_1 \right)\\
& &= 0,2 \cdot 0,001+0,3 \cdot 0,005+0,5 \cdot 0,006\\
& &= 0,0065
			\end{eqnarray*}
			\itemch Sai. Ta có $P\left(\overline{A}\right)=1-P\left(A\right)=1-0,0065=0,9935$.
			\itemch Sai. Ta có:$$P\left(A_3\vert A \right)=\dfrac{P\left(A_3\right)\cdot P\left(A \vert A_3\right)}{P\left(A\right)}=\dfrac{0,5\cdot 0,006}{0,0065}=\dfrac{6}{13}\approx 0,5$$
			Vậy xác suất để chiếc giày đó do nhà máy $3$ sản xuất là $0,5$.
		\end{itemchoice}
	}
\end{ex}

\Closesolutionfile{ans}
\indapan{3}{ans/ans-2-OTHK2}



\Opensolutionfile{ans}[ans/ans-2-OTHK2-De2-KQ]
\caukq
\begin{ex}%[Nguyễn Tuấn, dự án sáng tác đề 12]%[2H5H3-3]
	Trong không gian $Oxyz$, cho hai điểm $A(2;0;-1)$ và $B(0;2;3)$. Mặt cầu $(S)$ đường kính $AB$ có bán kính $R$. Khi đó $8R^2$ bằng
	\shortans{$48$}
	\loigiai{
		Gọi $I$ là trung điểm của $AB$.\\ 
		Khi đó, mặt cầu $(S)$ có tâm $I$ và bán kính $R=\dfrac{AB}{2}$.\\
		$\Rightarrow I(1;1;1)$, $R=\dfrac{1}{2} \sqrt{(0-2)^2+(2-0)^2+(3+1)^2}=\sqrt{6}$.\\
		Vậy $8R^2=48$.
	}
\end{ex}
\begin{ex}%[Nguyễn Tuấn, dự án sáng tác đề 12]%[2H5V2-2]
	Trong không gian $Oxyz$, cho điểm $M(1;0;2)$ và đường thẳng $d\colon \dfrac{x}{2}=\dfrac{y-1}{-1}=\dfrac{z-3}{1}$. Đường thẳng $\Delta$ đi qua $M$ cắt và vuông góc với $d$ có một véc-tơ chỉ phương là $\overrightarrow{u}=(a;b;4)$. Giá trị biểu thức $S=a+b$ bằng
	\shortans{$1$}
	\loigiai{
		Giả sử $H$ là giao điểm của hai đường thẳng $\Delta$ và $d$. Khi đó $H(2t;1-t;3+t)$.\\
		Vì $\Delta\perp d$ nên $\overrightarrow{MH}\cdot \overrightarrow{u}_d=0$.\\
		Ta có $\overrightarrow{MH}=(2t-1;1-t;1+t)$, $\overrightarrow{u}_d=(2;-1;1)$, suy ra
		$$2\cdot (2t-1)-1\cdot (1-t)+1\cdot (1+t)=0\Leftrightarrow t=\dfrac{1}{3}.$$
		Suy ra $\overrightarrow{MH}=\left(-\dfrac{1}{3};\dfrac{2}{3};\dfrac{4}{3}\right)$, vậy $\overrightarrow{u}=(-1;2;4)$ cũng là véc-tơ chỉ phương của $\Delta$.\\
		Vậy $a=-1$, $b=2$ hay $a+b=1$.
	}
\end{ex}
\begin{ex}%[Nguyễn Tuấn, dự án sáng tác đề 12]%[2D5V2-3] 
	Chuồng I có 5 con gà mái, 2 con gà trống. Chuồng II có 3 con gà mái, 5 con gà trống. Bác Mai bắt một con gà trong số đó theo cách sau: Bác tung một con xúc xắc cân đối, đồng chất. Nếu số chấm chia hết cho 3 thì bác chọn chuồng I, nếu số chấm không chia hết cho 3 thì bác chọn chuồng II. Sau đó, từ chuồng đã chọn bác bắt ngẫu nhiên một con gà. Gọi $P$ là xác suất để bác Mai bắt được con gà mái. Khi đó $84P$ bằng
	\shortans{$41$}
	\loigiai{
		Gọi $A$ là biến cố: \lq\lq Bác Mai bắt được con gà mái\rq\rq.\\
		Gọi $B_1$ là biến cố: \lq\lq tung con xúc xắc được số chấm chia hết cho $3$\rq\rq\, suy ra $\mathrm{P}(B_1)=\dfrac{2}{6}=\dfrac{1}{3}$.\\
		Gọi $B_2$ là biến cố: \lq\lq tung con xúc xắc được số chấm không chia hết cho $3$\rq\rq\, suy ra $\mathrm{P}(B_2)=\dfrac{4}{6}=\dfrac{2}{3}$.\\
		Ta có  xác suất bắt được gà mái từ chuồng I là $\mathrm{P}\left(A\mid B_1\right)=\dfrac{5}{7}$.\\
		Ta có  xác suất bắt được gà mái từ chuồng II là $\mathrm{P}\left(A\mid B_2\right)=\dfrac{3}{8}$.\\
		Áp dụng công thức xác suất toàn phần ta có
		$$\mathrm{P}(A)=\mathrm{P}(B_1)\cdot \mathrm{P}(A\mid B_1)+\mathrm{P}(B_2)\cdot \mathrm{P}(A\mid B_2)=\dfrac{1}{3}\cdot \dfrac{5}{7}+\dfrac{2}{3}\cdot \dfrac{3}{8}=\dfrac{41}{84}.$$
		Vậy xác suất để bác Mai bắt được con gà mái là $\dfrac{41}{84}$.\\
		Vậy $84P=41$.
	}
\end{ex}

\begin{ex}%[Nguyễn Tuấn, dự án sáng tác đề 12]%[2D4H3-3]
	\immini
	{
		Một vật trang trí có dạng khối tròn xoay tạo thành khi quay miền $(R)$ được giới hạn bởi đường gấp khúc $DABFE$ và cung tròn $ED$ (phần gạch chéo trong hình bên) xung quanh trục $AB$. Biết $ABCD$ là hình chữ nhật cạnh $AB =3$ cm, $AD=2$ cm; $F$ là trung điểm của $BC$; điểm $E$ cách $AD$ một đoạn bằng $1$ cm. Tính thể tích của vật trang trí đó, làm tròn kết quả đến hàng phần mười, đơn vị cm$^3$.
		\shortans{$16{,}5$}
	}
	{
		\begin{tikzpicture}[scale=1.2, line join=round, line cap=round,>=stealth,font=\footnotesize]
			\path (0,0) coordinate (B)	(2,0) coordinate (C) (2,3) coordinate (D) (0,3) coordinate (A) ($(B)!.5!(C)$) coordinate (F) ($(F)+(0,2)$) coordinate (E) ($(A)!.5!(D)$) coordinate (I) ;
			\draw (A)--(B)--(C)--(D)--cycle (E)--(F);
			\draw[dashed] (E)--(I);
			\foreach \d/\g in {A/90,B/-90,C/0,D/90,F/-90,E/-60}
			\draw[fill=black] (\d) circle (1pt) +(\g:0.2) node{$\d$};
			\draw (E) arc (-90:0:1cm); 
			\fill[color=gray,opacity=0.5] (A)--(B)--(F)--(I)--cycle (E) arc (-90:0:1cm)--(D)--(I)--cycle;
		\end{tikzpicture}
	}
	\loigiai{
		\immini
		{
			Gắn hệ trục tọa độ như hình vẽ bên.\\
			Quay hình chữ nhật $GEFB$ xung quanh trục $AB$ ta được khối trụ có bán kính đáy bằng $1$ và chiều cao bằng $2$.\\
			Phương trình đường tròn tâm $I$ bán kính bằng $1$ là $x^2+(y-1)^2=1$, suy ra phương trình của cung $DE$ là $y=1+\sqrt{1-x^2}$.\\
			Vây thể tích của vật trang trí là
			$$V=\pi \displaystyle\int\limits_{0}^{1} \left(1+\sqrt{1-x^2}\right)^2 \mathrm{\,d}x+\pi \cdot 1^2\cdot 2\approx 16{,}5\ (\mathrm{cm}^3).$$
		}
		{
			\begin{tikzpicture}[scale=1.2, line join=round, line cap=round,>=stealth,font=\footnotesize]
				\draw[->] (-0.5,0)--(4,0) node[below]{$x$};
				\draw[->] (0,-0.5)--(0,3) node[left]{$y$};
				\path (0,0) coordinate (A)	(3,0) coordinate (B) (3,2) coordinate (C) (0,2) coordinate (D) ($(B)!.5!(C)$) coordinate (F) ($(F)-(2,0)$) coordinate (E) ($(A)!.5!(D)$) coordinate (I) (1,0) coordinate (G) ;
				\draw (A)--(B)--(C)--(D)--cycle (E)--(F);
				\draw[dashed] (E)--(I) (E)--(G);
				\foreach \d/\g in {A/-120,B/-90,C/0,D/120,F/0,E/-60,I/180,G/-90}
				\draw[fill=black] (\d) circle (1pt) +(\g:0.2) node{$\d$};
				\draw (E) arc (0:90:1cm); 
				\fill[color=gray,opacity=0.5] (A)--(B)--(F)--(I)--cycle (E) arc (0:90:1cm)--(D)--(I)--cycle;
			\end{tikzpicture}	
		}
	}
\end{ex}
\begin{ex}%[Nguyễn Tuấn, dự án sáng tác đề 12]%[2D4C3-1]
	Cho hàm số bậc bốn $y=f(x)$ có ba điểm cực trị dương lần lượt là $x_1$, $x_2$, $x_3$ thỏa mãn $x_1+x_2+x_3=3$ và $g(x)$ là parabol đi qua ba điểm cực trị của đồ thị hàm số $f(x)$. Diện tích hình phẳng giới hạn bởi đồ thị của hàm số $y=\dfrac{f'(x)}{f(x)-g(x)}$, trục hoành và hai đường thẳng $x=-2$, $x=0$ bằng (kết quả làm tròn đến hàng phần trăm).
	\shortans{$4{,}39$}
	\loigiai{
		Giả sử $f(x)=ax^4+bx^3+cx^2+dx+e$ với $a$, $b$, $c$, $d$, $e\in\mathbb{R}$.\\
		Ta có $f'(x)=4ax^3+3bx^2+2cx+d=4a(x-x_1)(x-x_2)(x-x_3)$.\\
		Mà theo Định lý Vi-ét ta có $x_1+x_2+x_3=3\Leftrightarrow -\dfrac{3b}{4a}=3$.\\
		Vì $g(x)$ là hàm số bậc hai đi qua $3$ điểm cực trị của đồ thị $f(x)$ nên phương trình $f(x)=g(x)$ nhận $x_1$, $x_2$, $x_3$ là nghiệm. Vậy
		$$f(x)-g(x)=ax^4+bx^3+\left(cx^2+dx+e-g(x)\right)=a(x-x_1)(x-x_2)(x-x_3)(x-x_4).$$
		Ta có $x_1+x_2+x_3+x_4=-\dfrac{b}{a}=4$ mà $x_1+x_2+x_3=3$ nên $x_4=1$.\\
		Vậy $\dfrac{f'(x)}{f(x)-g(x)}=\dfrac{4}{x-1}$. Suy ra
		$$S=\displaystyle\int\limits_{-2}^{0} \left|\dfrac{4}{x-1}\right| \mathrm{\,d}x=4\ln 3\approx 4{,}39.$$
	}
\end{ex}
\begin{ex}%[Nguyễn Tuấn, dự án sáng tác đề 12]%[2H5C3-3]
	Trong không gian với hệ tọa độ $Oxyz$ cho điểm $A(-2;2;-2)$ và điểm $B(3;-3;3)$. Điểm $M$ thay đổi trong không gian thỏa mãn $\dfrac{MA}{MB}=\dfrac{2}{3}$. Điểm $N(a;b;c)$ thuộc mặt phẳng $(P)\colon-x+2y-2z+6=0$ sao cho $MN$ nhỏ nhất. Tính tổng $T=a+b+c$.
\shortans{$-2$}
	\loigiai{
		\immini
		{
			Gọi $M(x;y;z)$, ta có $\overrightarrow{MA}=(-2-x;2-y;-2-z)$,\\ $\overrightarrow{MB}=(3-x;-3-y;3-z)$. Vậy $$\dfrac{MA}{MB}=\dfrac{2}{3} \Leftrightarrow 9MA^2=4MB^2 \Leftrightarrow (x+6)^2+(y-6)^2+(z+6)^2=108.$$
			Vậy điểm $M$ thuộc mặt cầu tâm $I(-6;6 ;-6)$ bán kính $R=6\sqrt{3}$.\\
			Vậy $MN$ nhỏ nhất khi $M, N$ thuộc đường thẳng đi qua tâm $I$ và vuông góc với mặt phẳng $(P)$.\\
			Gọi $(d)$ là đường thẳng đi qua tâm $I$ và vuông góc với mặt phẳng $(P)$. Ta được $(d)\colon \heva{&x=-6-t \\& y=6+2t \\& z=-6-2t}$.\\
			Tọa độ điểm $N$ là nghiệm của hệ phương trình
			$$
			\begin{aligned}
				&\heva{& {x=-6-t} \\& {y=6+2t} \\& {z=-6-2t} \\& {-x+2y-2z+6=0}} \Leftrightarrow \heva{& {x=-6-t} \\& {y=6+2t} \\& {z=-6-2t} \\& {6+t+12+4t+12+4t+6=0}} \Leftrightarrow \heva{& x=-2 \\& y=-2 \\& z=2 \\& t=-4.}
			\end{aligned}
			$$
			Suy $N(-2;-2;2)$ nên ta được $T=-2-2+2=-2$.
		}
		{
			\begin{tikzpicture}[scale=.6, font=\footnotesize, line join=round, line cap=round, >=stealth]
				\def\R{2.5} 
				\def\t{.15} 
				\coordinate[label=left:$I$] (O) at (0,0);
				\coordinate (A) at (\R,0);
				\draw (O) circle (\R) (A) arc (0:-180: {\R} and {\R/3});
				\draw[dashed] (O)--(A) arc (0:180: {\R} and {\R/3});
				\foreach \diem in {A,O}	\fill (\diem)circle(1pt);
				\coordinate (M) at ($(0,0)+(-45:\R cm)$);
				\coordinate (N) at ($(0,0)+(-135:\R cm)$);
				\draw[dashed] (M)--(N);
				\draw (M)--++(0:1.5)--++(-135:2.5)coordinate (I)--++(180:5.5)coordinate (J) --++(45:2.5)coordinate (K)--(N);
				\draw[dashed] (O)--++(-90:\R)--++(-90:0.5)coordinate (H);
				\draw (H)--++(90:0.5);
				\fill (H) circle (1pt) node[below]{$H$};
				\def\t{.15}
				\draw (H)--++(90:\t)--++(0:\t)--++(-90:\t);
				\clip (I) -- (J) -- (K);
				\draw (J) circle (1cm);
				\draw ($(J)+(0.25,0)$) node[above right]{$(P)$};
			\end{tikzpicture}
			
		}
		
	}
\end{ex}
\Closesolutionfile{ans}
\indapan{6}{ans/ans-2-OTHK2-De2-KQ}
